Bên cạnh các yêu cầu về độ chính xác, độ tin cậy hay khả năng chống chịu của hệ thống, tác động xã hội (Social Impact) là một trong những khía cạnh quan trọng của khung Trustworthy AI. Tiêu chí này tập trung đánh giá ảnh hưởng của hệ thống trí tuệ nhân tạo đối với con người và cộng đồng, đặc biệt là các nhóm người dùng yếu thế có nguy cơ chịu tác động bất lợi từ công nghệ \cite{liu2021_trustworthyai}.

Đối với đề tài chatbot phát hiện lừa đảo điện thoại cho người cao tuổi, việc đánh giá tác động xã hội mang ý nghĩa đặc biệt quan trọng bởi đối tượng sử dụng chính là nhóm người có mức độ dễ bị tổn thương cao trước các hình thức lừa đảo công nghệ. Nếu được thiết kế phù hợp, hệ thống có thể góp phần nâng cao khả năng tự bảo vệ, giảm thiểu thiệt hại tài chính và tăng cường an toàn số cho người cao tuổi. Ngược lại, nếu không được xem xét đầy đủ dưới góc độ xã hội và đạo đức, chatbot có thể tạo ra những hệ quả ngoài mong muốn như làm gia tăng sự phụ thuộc vào AI, tạo khoảng cách số hoặc ảnh hưởng đến quyền riêng tư và quyền tự chủ của người sử dụng.
Do đó, việc phân tích cả tác động tích cực và các rủi ro tiềm ẩn là cần thiết nhằm đảm bảo hệ thống mang lại lợi ích thực sự cho cộng đồng và phù hợp với nhu cầu của nhóm người dùng mục tiêu.

\section{Tác động tích cực}
\subsection{Hỗ trợ bảo vệ tài chính cho người cao tuổi}
Một trong những lợi ích quan trọng nhất của chatbot là khả năng hỗ trợ người cao tuổi phòng tránh các hình thức lừa đảo tài chính qua điện thoại. Trong thực tế, phần lớn người cao tuổi phụ thuộc vào lương hưu hoặc các khoản tiết kiệm được tích lũy trong nhiều năm. Khi trở thành nạn nhân của các vụ lừa đảo, họ thường phải đối mặt với những tổn thất tài chính nghiêm trọng và khó có khả năng khắc phục trong thời gian ngắn.

Thông qua việc phân tích nội dung cuộc gọi và phát hiện các dấu hiệu bất thường như yêu cầu chuyển tiền, cung cấp mã OTP hoặc tạo áp lực tâm lý, chatbot có thể đưa ra cảnh báo sớm giúp người dùng nhận biết nguy cơ và tránh đưa ra các quyết định thiếu an toàn. Điều này góp phần giảm thiểu nguy cơ mất tài sản và hạn chế các hậu quả kinh tế đối với cá nhân cũng như gia đình người bị hại.

\subsection{Giảm thiểu tổn thương tâm lý}
Bên cạnh thiệt hại tài chính, các vụ lừa đảo còn để lại nhiều hậu quả về mặt tâm lý đối với nạn nhân. Nhiều người cao tuổi sau khi bị lừa đảo thường trải qua cảm giác tự trách bản thân, mất niềm tin vào người khác hoặc lo sợ khi tiếp nhận các cuộc gọi từ số lạ. Những tác động này có thể trở nên nghiêm trọng hơn đối với người sống một mình hoặc thiếu sự hỗ trợ từ gia đình.

Việc chatbot đưa ra cảnh báo kịp thời và hướng dẫn các bước xác minh thông tin có thể giúp người dùng cảm thấy an tâm hơn trong quá trình sử dụng điện thoại và các dịch vụ số. Qua đó, hệ thống góp phần giảm thiểu áp lực tâm lý, tăng cảm giác an toàn và hạn chế khả năng bị thao túng trong các tình huống khẩn cấp.

\subsection{Tăng cường khả năng tự chủ trong môi trường số}
Một lợi ích xã hội đáng chú ý khác là khả năng hỗ trợ người cao tuổi duy trì sự độc lập trong quá trình sử dụng công nghệ. Nhiều người lớn tuổi hiện nay thường phải nhờ con cháu hoặc người thân hỗ trợ khi gặp các cuộc gọi đáng ngờ hoặc các tình huống liên quan đến giao dịch trực tuyến.

Sự xuất hiện của chatbot giúp người dùng có thêm một công cụ hỗ trợ để tự đánh giá mức độ rủi ro của cuộc gọi, từ đó chủ động hơn trong việc ra quyết định. Điều này không chỉ giúp nâng cao kỹ năng an toàn số mà còn góp phần duy trì cảm giác tự chủ và phẩm giá cá nhân của người cao tuổi trong xã hội số.

\subsection{Nâng cao nhận thức cộng đồng về an toàn số}
Ở phạm vi rộng hơn, việc triển khai các công cụ hỗ trợ chống lừa đảo có thể mang lại những lợi ích tích cực cho toàn xã hội. Các cảnh báo và giải thích do chatbot cung cấp giúp người dùng hiểu rõ hơn về các hình thức lừa đảo phổ biến, từ đó nâng cao nhận thức về an toàn thông tin và kỹ năng phòng tránh rủi ro trên môi trường số.

Ngoài ra, dữ liệu tổng hợp và được ẩn danh từ hệ thống có thể hỗ trợ việc nhận diện các xu hướng lừa đảo mới, góp phần nâng cao hiệu quả của các chương trình tuyên truyền và bảo vệ người dân trước các mối đe dọa công nghệ trong tương lai.

\section{Tác động tiêu cực và rủi ro}
\subsection{Nguy cơ phụ thuộc quá mức vào AI}
Một trong những rủi ro đáng chú ý là người dùng có thể dần phụ thuộc hoàn toàn vào chatbot trong quá trình đánh giá các cuộc gọi đáng ngờ. Nếu điều này xảy ra, người cao tuổi có thể giảm khả năng tự xác minh thông tin hoặc không còn duy trì thói quen tham khảo ý kiến từ người thân, ngân hàng hoặc cơ quan chức năng.

Trong trường hợp hệ thống gặp lỗi hoặc không thể hoạt động, sự phụ thuộc quá mức này có thể khiến người dùng dễ trở thành nạn nhân của các hình thức lừa đảo mà không có cơ chế bảo vệ thay thế. Hậu quả nghiêm trọng nhất là chatbot đánh giá nhầm một cuộc gọi lừa đảo là an toàn. Khi đó, người dùng có thể đặt niềm tin tuyệt đối vào hệ thống và thực hiện các hành động nguy hiểm như chuyển tiền hoặc cung cấp thông tin cá nhân.
Đối với người cao tuổi, rủi ro này đặc biệt nghiêm trọng bởi họ thường có xu hướng tin tưởng vào các hệ thống công nghệ được giới thiệu như một công cụ hỗ trợ chính thức. Vì vậy, việc đưa ra các kết luận mang tính khẳng định tuyệt đối có thể dẫn đến những hậu quả khó lường.

Khi người dùng quá phụ thuộc vào AI, một hệ thống có tỷ lệ cảnh báo sai cao có thể khiến người dùng trở nên lo lắng hoặc nghi ngờ ngay cả đối với những cuộc gọi hoàn toàn hợp pháp. Về lâu dài, điều này có thể làm giảm niềm tin vào công nghệ hoặc khiến người dùng bỏ qua các cảnh báo thực sự quan trọng do đã quá quen với các cảnh báo giả.
Đối với người cao tuổi, tác động tâm lý từ các cảnh báo sai thường lớn hơn do họ có xu hướng nhạy cảm hơn với các thông tin liên quan đến an toàn cá nhân và tài chính.

\subsection{Nguy cơ làm suy giảm quyền tự chủ của người dùng}
Một số giải pháp hỗ trợ người cao tuổi có thể được thiết kế theo hướng tự động gửi cảnh báo cho người thân hoặc can thiệp vào quá trình liên lạc của người dùng. Mặc dù mục tiêu của các tính năng này là bảo vệ người sử dụng, chúng cũng có thể dẫn đến việc hạn chế quyền tự chủ cá nhân.
Nếu không được triển khai một cách minh bạch và có sự đồng thuận rõ ràng từ người dùng, hệ thống có thể khiến người cao tuổi cảm thấy bị giám sát hoặc mất quyền kiểm soát đối với các quyết định cá nhân của mình.

\subsection{Gia tăng khoảng cách số}
Mặc dù chatbot được thiết kế nhằm hỗ trợ người cao tuổi, không phải tất cả người dùng trong nhóm này đều có điều kiện tiếp cận công nghệ như nhau. Những người sinh sống ở vùng sâu vùng xa, người có thu nhập thấp hoặc người có kỹ năng số hạn chế có thể gặp khó khăn trong việc sử dụng các ứng dụng dựa trên điện thoại thông minh hoặc Internet.
Nếu hệ thống chỉ được triển khai trên các nền tảng hiện đại, một bộ phận người dùng dễ bị tổn thương nhất có thể không được hưởng lợi từ công nghệ. Điều này có nguy cơ làm gia tăng khoảng cách số giữa các nhóm dân cư khác nhau.

\section{Giải pháp giảm thiểu tác động tiêu cực}
Để đảm bảo chatbot mang lại giá trị xã hội tích cực và hạn chế các rủi ro nêu trên, hệ thống cần được thiết kế theo nguyên tắc lấy con người làm trung tâm.

Trước hết, chatbot cần được định vị là công cụ hỗ trợ ra quyết định thay vì công cụ đưa ra quyết định thay cho người dùng. Các phản hồi của hệ thống nên mang tính khuyến nghị và luôn khuyến khích người dùng xác minh thêm thông tin thông qua người thân hoặc các nguồn chính thức.

Đối với vấn đề khoảng cách số, hệ thống cần hỗ trợ nhiều hình thức tiếp cận khác nhau thay vì chỉ giới hạn trên điện thoại thông minh. Các kênh như tin nhắn SMS, tổng đài tự động hoặc trợ lý giọng nói có thể giúp mở rộng khả năng tiếp cận đến nhiều nhóm người dùng hơn.

Cuối cùng, giao diện và cơ chế tương tác cần được thiết kế phù hợp với đặc điểm của người cao tuổi. Các cảnh báo nên được trình bày bằng ngôn ngữ đơn giản, dễ hiểu và mang tính hướng dẫn thay vì tạo cảm giác hoảng sợ. Điều này giúp người dùng tiếp cận thông tin một cách hiệu quả hơn và giảm thiểu các tác động tiêu cực về mặt tâm lý.

\subsection{Thiết kế giao diện thân thiện với người cao tuổi}
Giao diện web của chatbot được thiết kế để người cao tuổi dễ đọc và dễ thao tác. Chữ trong tin nhắn và vùng nhập liệu được đặt ở cỡ lớn, rộng rãi hơn nhiều so với các ứng dụng chat thông thường, giúp người có thị lực kém vẫn đọc được. Khoảng cách giữa các dòng và lề trong khung tin nhắn cũng được nới rộng để nội dung không bị dồn nén. Các nút bấm có kích thước lớn, đủ chỗ để người dùng chạm chính xác ngay cả khi tay run hoặc mắt không nhìn rõ.

Màu sắc được chọn sao cho chữ nổi rõ trên nền trắng, đường viền dày giúp phân biệt các vùng chức năng. Nút hành động có màu xanh làm chủ đạo và hiệu ứng chuyển màu nhẹ để dễ nhận biết. Trạng thái kết nối được hiển thị bằng chấm tròn màu xanh kèm dòng chữ "Kết nối an toàn" để tạo cảm giác tin tưởng.

Thanh bên chứa lịch sử trò chuyện có thể thu gọn trên điện thoại, thay thế bằng nút mở menu và lớp phủ mờ. Nút bắt đầu cuộc trò chuyện mới (hình dấu cộng màu xanh lá) xuất hiện ở cả thanh bên lẫn đầu trang, luôn trong tầm với trên mọi kích thước màn hình. Các nút gợi ý chủ đề như "Có số lạ gọi cho tôi" giúp người dùng bắt đầu mà không cần gõ bàn phím.

Về hỗ trợ giọng nói, hệ thống có thể nhận diện giọng nói tiếng Việt và đọc văn bản thành tiếng. Khi người dùng bấm nút micro, dòng chữ hướng dẫn chuyển thành "Đang nghe ông bà nói..." để xác nhận hệ thống đang lắng nghe. Sau khi chatbot trả lời, nội dung được đọc tự động, giúp người khiếm thị hoặc không quen đọc chữ vẫn có thể sử dụng.

Khi nhận diện giọng nói thất bại, hệ thống hiển thị và đọc lên thông báo "Dạ cháu chưa nghe rõ, ông bà bấm lại micro để nói nhé ạ!" thay vì thông báo lỗi kỹ thuật. Khi mất kết nối, thông báo "Dạ hệ thống của cháu đang bận, ông bà thử lại sau nhé ạ!" xuất hiện, giữ đúng giọng điệu lịch sự và cách xưng hô.

Lịch sử trò chuyện được lưu trong bộ nhớ máy người dùng, không gửi lên máy chủ, đảm bảo quyền riêng tư. Người dùng có thể xóa từng cuộc trò chuyện hoặc tải toàn bộ lịch sử về dưới dạng tệp văn bản qua nút xuất ở cuối thanh bên.

Thông qua các biện pháp trên, chatbot có thể phát huy tối đa các lợi ích xã hội đồng thời hạn chế những rủi ro tiềm ẩn đối với các nhóm người dùng yếu thế, từ đó đáp ứng tốt yêu cầu về Social Impact trong khung đánh giá Trustworthy AI.

\section{Tiêu chí đánh giá}
Tác động xã hội của hệ thống được đánh giá thông qua:
\begin{itemize}
    \item Mức độ hỗ trợ người cao tuổi nhận diện lừa đảo.
\item Mức độ cải thiện nhận thức về an toàn số.
\item Tác động đến cảm giác an toàn và tự chủ của người sử dụng.
\item Nguy cơ phụ thuộc quá mức vào AI.
\item Mức độ chấp nhận của người dùng đối với hệ thống.
\end{itemize}
